\section{Vepurisikā, Dvipuruṣikā}
\label{vepurisika}


There word \textit{vepurisikā} in the Pali Canon does not appear anywhere other than in the Pali \textit{Bhikkhunikkhandhaka} section on the questons for ordination of a \textit{Bhikkhunī}. 

There is a similar word mentioned in the Sanskrit Canon that seems to be closely related and also appears in the same place in the section on the \textit{Bhikkhunī} ordination questions. The word \textit{dvipuruṣikā} appears in the only surviving copy on the \textit{Bhikṣuṇī Vinaya} of the Mahāsāṃghika-Lokottaravāda school\footnote{\cite{roth}. The Chinese pilgrim Faxian, who translated a large number of texts into Chinese, procured his physical Mahāsaṅghika manuscripts at a monastery in Pāṭaliputra. Pāṭaliputra was the historical cradle of the core, ancestral Mahāsaṅghika order. The Lokottaravāda branch branched off and flourished primarily in the northwest, particularly around Bamiyan (modern-day Afghanistan). Although Faxian did walk through northwestern India on his outward journey, he explicitly laments in his travelogue that in Northern and Northwestern India, the Vinaya was taught entirely through oral recitation and memory—there were no physical manuscripts available for him to purchase or copy (\cite{legge1886} pp. 98--99). When Xuanzang visited the Bamiyan Valley in 643 CE, he collected a lot of the Sanskrit texts of this school and brought them back to China. Unfortunately, all the original Sanskrit texts, mostly written on berch-bark, palmost all the physical Sanskrit manuscripts carried across Central Asia or the maritime Silk Roads by ancient Chinese pilgrims over the centuries.}. This word does not appear in any other known Sanskrit texts.

The two traditions approach the concept through slightly different semantic angles: 
\begin{table}[htbp]
\centering
\small
\renewcommand{\arraystretch}{1.3}
\begin{tabularx}{\textwidth}{@{}l l >{\raggedright\arraybackslash}X >{\raggedright\arraybackslash}X@{}}
\toprule
\textbf{Tradition} & \textbf{Term} & \textbf{Etymology} & \textbf{Literal meaning} \\
\midrule
Theravāda & \textit{vepurisikā} &
  \textit{vi-} (anomalous/deviant) + \textit{purisa} (man) + \textit{-ikā} (feminine suffix) &
  a deviant or anomalous man-woman \\
Lokottaravāda & \textit{dvipuruṣikā} &
  \textit{dvi-} (two) + \textit{puruṣa} (man/person) + \textit{-kā} (feminine suffix) &
  having the nature of two persons \\
\bottomrule
\end{tabularx}
\end{table}

The Theravāda commentary, the \textit{Samantapāsādikā}, glosses this term as: ''A  woman who has a beard and mustache, and possesses the physical form of a man''\footnote{Sp.1.285}.

The Lokottaravāda school, unlike the Theravāda school, has no surviving commentarial texts on which we can draw. From the etymological context, it seems to approach this concept as a phenotypic variant—specifically, a female person who has developed secondary male characteristics (such as a beard, mustache, deep voice, and/or heavily masculine frame).

There is only one other known surviving copy of a Sanskrit text that contains the ordination questions for \textit{Bhikkhunīs}, being of the Mūlasarvāstivāda school\footnote{\cite{schmidt}}. Two more copies of the ordination questions for \textit{Bhikkhus} survived of the Sarvāstivāda school\footnote{\cite{haertel1956} and \cite{chung2004}}. All these three texts do not mention the word \textit{dvipuruṣikā}, but they mention the word \textit{ubhayavyañjanā}, while the text from the Lokottaravāda school lists the \textit{dvipuruṣikā} only. 

It is also noteworthy in this case that the Theravāda school is the only school that mentions both the \textit{ubhatobyañjanā} (Sanskrit \textit{ubhayavyañjanā}) and the \textit{vepurisikā}, while all other schools only list one of them. Whenever there is a discrepancy where only one school lis(Sanskrit \textit{ubhayavyañjanā}) ts something where the others remain silent, it likely denotes a later expansion of the texts. 

In \cite{vimala2021}\footnote{\cite{vimala2021} pp. 26--31} I analyzed the term \textit{ubhatobyañjanā}. Although the word itself originated in context of the Buddhist Vinaya, probably in order to achieve the
Classical Indian (and especially Buddhist) fourfold logical tetralemma called the catuṣkoṭi\footnote{\cite{scherer} p. 68 notes that
the categories of \textit{paṇḍaka} and \textit{ubhatobyañjanā} are largely academic and intended to make up the fourfold taxonomy (‘male’, ‘female’, ‘both ...’, ‘neither ...’).}, the whole concept was also bound up with older mythological stories of shape-shifting hermaphroditic beings who can reproduce as both a male and a female as they have a full set of reproductive organs of each.

Considering the close proximity in meaning between the \textit{dvipuruṣikā} and \textit{ubhayavyañjanā}, it might be possible that the same concept is meant. Also the fact that only two schools make any mention of this concept, while all the Chinese translations, of which the Sanskrit originals were lost, do not, warrents a closer look at the Chinese texts.

In \cite{vimala2021}\footnote{\cite{vimala2021} pp. 16--26} I traced the Pali term \textit{paṇḍaka} back to ancient Vedic mythology and social-religious ritual classes. This culture was very different from the Chinese and because the Chinese translators did not have an adequate word for this term, they used concepts from their own culture. It would not be surprising if a similar textual confusion has happened with the word \textit{dvipuruṣikā}. 

In Chinese Buddhist translations and Abhidharma psychology, two distinct terms are often used when navigating sex- and gender-characteristics. While sometimes used as synonyms\footnote{When later Chinese masters, such as Daoxuan (道宣) of the Tang Dynasty, wrote authoritative commentaries on monastic law (such as the Sifen-lu Shanfan Bujue Jiemo, 隨機羯磨), they synthesized these terms together. It is in these Chinese commentaries—rather than the original Taisho Mahāsāṅghika text—that we will find an explicit statement that 二形 means a body possessing 二根 (both faculties):  T40n1808_p0497b13 汝非二形耶？（謂此身中具有男女根，正乖道器。 Translation: ''Are you not a 二形? (This refers to one who possesses both the male and female faculties within this body; such a person is fundamentally unsuitable as a vessel for the Way.)''}, they maintain a strict technical separation:

\begin{table}[htbp]
\centering
\small
\renewcommand{\arraystretch}{1.3}
\begin{tabularx}{\textwidth}{@{}l l >{\raggedright\arraybackslash}X@{}}
\toprule 
\textbf{Term} & \textbf{Literal meaning} & \textbf{Gloss} \\
\midrule
二形 (èr xíng) & dual form / signs &
  Refers to external morphology and secondary sex characteristics. \\
二根 ((èr gēn) & dual faculties / organs &
  Refers strictly to anatomy. It requires the concurrent physical presence of both functional sets of genitalia. \\
\bottomrule
\end{tabularx}
\end{table}

It seems from the above that 二形 is more akin to \textit{vepurisikā/dvipuruṣikā} while 二根 refers more to \textit{ubhatobyañjanā/ubhayavyañjanā}. But are these terms used in the same way throughout the texts?

As I mentioned above, there are three surviving versions of the ordination questions from two schools of which we know there are versions translated into Chinese. All three mention the word \textit{ubhayavyañjanā}. 


- ---------

The 7th-century CE Kāśikāvṛtti\footnote{The Kāśikāvṛtti is a running commentary (vṛtti) on the grammatical aphorisms of Pāṇini's Aṣṭādhyāyī, jointly authored by Jayāditya and Vāmana. }, mentions:



This exact passage and its specific phrasing were quoted directly from Patañjali’s Mahābhāṣya (c. 150 BCE), who in turn was debating a verses-based biological theory that dates all the way back to the Vedic and early post-Vedic anatomical schools (c. 500 BCE).
The specific history of how this anatomical definition ended up in a Sanskrit grammar text follows a clear evolutionary line:

1. The Pre-Grammatical Origin: Early Indian Medical and Structural Thought (c. 500 BCE)

Before it became a tool for grammar, the distinction was a basic rule of thumb in ancient Indian observational science. To classify people visible in the everyday world (loka), early texts identified clear visual markers:
• Feminine marker: The presence of breasts (stana).
• Masculine marker: The growth of facial and body hair (loma / śmaśru).
• The Dual State: A person who physically developed both distinct markers simultaneously was classified as having "both signs" (ubhayavyañjana).

2. Patañjali’s Structural Leap in the Mahābhāṣya (c. 150 BCE)

https://gretil.sub.uni-goettingen.de/gretil/1_sanskr/6_sastra/1_gram/pmbhasuu.htm

The exact verse structure you see preserved in the Kāśikā was popularized by Patañjali in his commentary on Pāṇini's sūtra 4.1.3 (which deals with feminine suffixes). Patañjali quotes older traditional verses (ślokas) to test if grammar can simply rely on physical appearance to assign grammatical gender:
स्तनकेशवती स्त्री स्याल्लोमशः पुरुषः स्मृतः । उभयोर्न्तरं यच्च तदभावे नपुंसकम् ॥
"She who has breasts and long hair is a woman; he who is hairy is a man. That which is between the two, or lacks both, is a napuṃsaka."
Patañjali brings up the ubhayavyañjana concept to point out a fundamental flaw: if gender in grammar only meant physical features, then a table (khaṭvā) wouldn't be feminine because it doesn't have breasts, and an intersex person with both physical traits would break the rules of language. Patañjali introduces the somatic definition specifically to reject it as a complete explanation for grammar, choosing instead a more abstract theory of gender based on the primary forces of emergence (āvirbhāva) and disappearance (tirobhāva).





